% Appendix (formal stationary phase), part 1 of 5; label prefix: fsp. Each part inputs the next one at its end.
% No new macros; local letters are declared where they occur.
\section{Formal stationary phase}
\label{sec:fsp}

This appendix proves the properties of formal stationary phase used in \cref{sec:con,sec:gam}. Coefficients lie in a commutative $\bbQ$-algebra $A$, the stationary phase at a critical point is the Wick contraction of the Taylor data there, and nothing is evaluated at $z\neq0$. Only the branch part of \cref{lem:fsp-refl}, used in \cref{lem:con-form}~(d), is not about formal series. Here $t,x,y,v,w$ denote integration or outer variables, $\xi$ the variable of a generating function, $q$ and $\varepsilon$ formal parameters, and $\mathsf H$ a Hessian matrix.

%-------------------------------------------------------------------------------
\subsection{The weighted ring and the stationary-phase functional}
\label{sec:fsp-def}

\begin{defn}[Weighted ring and Gaussian functional]
\label{def:fsp-ring}
Let $t=(t_1,\dots,t_n)$ be variables and $z$ a formal variable, and put $\wt(t^\alpha z^j):=|\alpha|/2+j$ for $\alpha\in\bbZ_{\geq0}^n$, $j\in\bbZ$. Let $\cR_n(A)$ be the set of formal sums $\sum c_{\alpha,j}t^\alpha z^j$, $c_{\alpha,j}\in A$, such that for every $W$ only finitely many monomials of weight $\leq W$ have $c_{\alpha,j}\neq0$. A family $(P_\lambda)$ in $\cR_n(A)$ is \emph{summable} if for every $W$ only finitely many $P_\lambda$ have a non-zero term of weight $\leq W$; its sum is taken termwise. For a symmetric $C\in A^{n\times n}$ put $\langle t^\alpha\rangle_C:=\de_\xi^\alpha\exp(\xi^{\mathrm T}C\xi/2)|_{\xi=0}$, so that $\sum_\alpha\langle t^\alpha\rangle_C\,\xi^\alpha/\alpha!=\exp(\xi^{\mathrm T}C\xi/2)$, extended $A$-linearly to $A[t]$. The \emph{Gaussian functional} with covariance $zC$ is the $A$-linear map $E_{zC}\colon\cR_n(A)\to A\lser{z}$, $E_{zC}[t^\alpha z^j]:=\langle t^\alpha\rangle_C\,z^{j+|\alpha|/2}$, applied termwise; $E_{zC}[1]=1$.
\emph{Outer variables.} If variables $y=(y_1,\dots,y_m)$ of weight $1/2$ are added and $C$ has entries in $A\fps{y}$, the partial functional $E^t_{zC}\colon\cR_{n+m}(A)\to\cR_m(A)$, $t^\alpha y^\beta z^j\mapsto\langle t^\alpha\rangle_C\,y^\beta z^{j+|\alpha|/2}$, acts on $t$ only.
\emph{Parameters.} Over $A=A_0\fps{q}$ we work in $\cR_n(A_0)\fps{q}\supset\cR_n(A)$, coefficientwise in $q$, in the following sense. For a symmetric matrix $C$ over $A$, which may depend on $q$, and $P=\sum_dq^dP_d$ with $P_d\in\cR_n(A_0)$, put $E_{zC}[P]:=\sum_dq^dE_{zC}[P_d]\in A_0\lser z\fps q$, where $E_{zC}[P_d]\in A\lser z$; on $\cR_n(A)$ this is the functional defined above. Modulo $q^N$ it is the Gaussian functional over $B_N:=A_0[q]/(q^N)$ on $\cR_n(B_N)=\cR_n(A_0)\fps q/(q^N)$. Identities in $\cR_n(A_0)\fps q$, in particular between exponentials, are meant modulo every $q^N$, in $\cR_n(B_N)$, where an element divisible by $q$ is nilpotent and $\re^{Y+X}=\re^Y\re^X$ for $Y$ of weight $\geq1/2$ and $X$ nilpotent. Likewise, for an amplitude $a\in A_0\fps{x-c_0}\lser z\fps q$ and a critical point $c\in c_0+(qA)^n$, the series $\SPn_c[F;a]\in A_0\lser z\fps q$ of \cref{def:fsp-wick} is defined modulo every $q^N$, where $a(c+t)$ is an amplitude over $B_N$. Alternatively $q$ gets the weight $1$, and $\cR_{n;q}(A_0)$ is the set of sums $\sum c\,t^\alpha q^dz^j$ with the same finiteness condition for the weight $|\alpha|/2+d+j$.
\end{defn}

\begin{lem}[The weighted ring]\genusstatus{AI-proposed, unchecked}
\label{lem:fsp-ring}
(i) $\cR_n(A)$ is a ring containing $A\fps{t}\lser{z}$; multiplication by a fixed element, $\de/\de t_i$ and multiplication by $t_i$ preserve summable families and commute with their sums.

(ii) If every term of $P\in\cR_n(A)$ has weight $\geq1/2$, the family $(P^k/k!)_k$ is summable, $\exp(P):=\sum_kP^k/k!$ lies in $\cR_n(A)$, and the usual rules for exponentials hold.

(iii) $\langle t^\alpha\rangle_C=0$ for $|\alpha|$ odd; for $|\alpha|$ even it is the sum over the perfect matchings of $|\alpha|$ labelled points, $\alpha_i$ of them carrying the index $i$, of the products of the $C_{ij}$ over the matched pairs; and $\langle t_it^\beta\rangle_C=\sum_jC_{ij}\beta_j\langle t^{\beta-e_j}\rangle_C$ (Wick recursion).

(iv) $E_{zC}$ and $E^t_{zC}$ are well defined, do not lower the weight, and commute with sums of summable families; $E_{zC}$, and $E^t_{zC}$ for constant $C$, preserve the weight, $E_{zC}[fP]=fE_{zC}[P]$ for $f\in A\lser{z}$, and $E^t_{zC}[gP]=g\,E^t_{zC}[P]$ for $g\in\cR_m(A)$, a series in the outer variables and $z$, and $P\in\cR_{n+m}(A)$.

(v) For $J\in A^{n\times n}$ the substitution $t=Jy$ maps $\cR_n(A)$ to itself preserving the weight, and $E_{zC'}[P(Jy)]=E_{zJC'J^{\mathrm T}}[P]$ for all $P\in\cR_n(A)$ and symmetric $C'$.
\end{lem}

\begin{proof}
(i) The weights of the non-zero terms of an element are bounded below, so in a product the coefficient of a monomial of weight $W$ is a finite sum over pairs of terms; the same count gives the finiteness condition and the assertions on summable families. An element $\sum_{j\geq j_0}a_jz^j$ of $A\fps{t}\lser{z}$ has in weight $\leq W$ only terms with $j_0\leq j\leq W$, $|\alpha|\leq2(W-j_0)$. The operators $\de/\de t_i$ and $t_i\cdot$ shift weights by $\mp1/2$. (ii) $P^k$ has weight $\geq k/2$. (iii) $\exp(\xi^{\mathrm T}C\xi/2)$ is even in $\xi$, and $\de_{\xi_i}\exp(\xi^{\mathrm T}C\xi/2)=(C\xi)_i\exp(\xi^{\mathrm T}C\xi/2)$; applying $\de_\xi^\beta$ at $\xi=0$ with the Leibniz rule gives the recursion (one derivative $\de_{\xi_j}$ must hit $\xi_j$, in $\beta_j$ ways). The recursion determines all moments from $\langle1\rangle_C=1$, and the matching sum satisfies it too (match $i$ with one of the remaining entries). (iv) A monomial of weight $w$ goes to a multiple of $z^w$, odd moments vanishing, and the coefficient of $z^k$ receives contributions only from the finitely many terms of weight $k$; this gives the assertions on $E_{zC}$. For $E^t_{zC}$, $\langle t^\alpha\rangle_C\in A\fps{y}$ is a polynomial in the entries of $C$, so the image of $t^\alpha y^\beta z^j$ is a sum of monomials $y^{\beta+\gamma}z^{j+|\alpha|/2}$ of weight $\geq\wt(t^\alpha y^\beta z^j)$ ($\gamma=0$ if $C$ is constant), and the coefficient of $y^\delta z^k$ receives contributions only from the finitely many terms with $\beta\leq\delta$, $j+|\alpha|/2=k$, which have weight $\leq k+|\delta|/2$. For $g\in\cR_m(A)$ let $g_\gamma$ be its monomials; the family $(g_\gamma P)_\gamma$ is summable with sum $gP$, and $E^t_{zC}[g_\gamma P]=g_\gamma E^t_{zC}[P]$ by the definition on monomials. (v) By the generating function, $\langle\re^{\xi\cdot Jy}\rangle_{C'}=\langle\re^{(J^{\mathrm T}\xi)\cdot y}\rangle_{C'}=\exp(\xi^{\mathrm T}JC'J^{\mathrm T}\xi/2)$, i.e.\ $\langle(Jy)^\alpha\rangle_{C'}=\langle t^\alpha\rangle_{JC'J^{\mathrm T}}$, which is the claim for $P=t^\alpha z^j$; the substitution maps a monomial of weight $w$ to a finite sum of monomials of weight $w$, and (iv) extends the identity.
\end{proof}

\begin{defn}[The normalised stationary phase]
\label{def:fsp-wick}
Let $c\in A^n$. For a \emph{phase} $F\in A\fps{x-c}$, $c$ is a \emph{non-degenerate critical point over $A$} if $\nabla F(c)=0$ and $\mathsf H:=\Hess F(c)\in GL_n(A)$. With $t:=x-c$ write $F(c+t)=F(c)+\frac12t^{\mathrm T}\mathsf Ht+F_{\geq3}(t)$, $F_{\geq3}\in(t)^3$, and $C:=-\mathsf H^{-1}$. An \emph{amplitude} is a series $a\in A\fps{x-c}\lser{z}$, or any $a$ with $\exp(F_{\geq3}/z)\,a(c+t)\in\cR_n(A)$; every term of $F_{\geq3}/z$ has weight $\geq1/2$. The \emph{normalised stationary phase} is
\[
\SPn_c[F;a]:=E_{zC}\bigl[\exp\bigl(F_{\geq3}(t)/z\bigr)\,a(c+t)\bigr]\in A\lser{z}.
\]
Equivalently, substitute $t=z^{1/2}\tilde t$, expand in powers of $z^{1/2}$ and apply $\langle\cdot\rangle_C$ in $\tilde t$ termwise. The map $a\mapsto\SPn_c[F;a]$ is $A\lser{z}$-linear, in particular $A\fps{z}$-linear (\cref{lem:fsp-ring}~(iv)). If $a=\sum_{m\geq0}a_mz^m$, the constant term of $\SPn_c[F;a]$ is $a_0(c)$; if it is a unit, $\SPh_c[F;a]:=\SPn_c[F;a]/a_0(c)\in1+zA\fps{z}$. The \emph{full symbol} is
\[
\mathrm{SP}_c\bigl[\re^{F/z}a\,\dd^nx\bigr]:=\bigl[\re^{F(c)/z}\bigr]\,(2\pi z)^{n/2}\,\bigl[\det(-\mathsf H)^{-1/2}\bigr]\,\SPn_c[F;a],
\]
with formal symbols in brackets ($[\re^{F_1/z}][\re^{F_2/z}]=[\re^{(F_1+F_2)/z}]$ for $F_1,F_2\in A$; the square root, in an extension of $A$, is part of the datum). An identity of full symbols is an identity of normalised series together with an explicit identity of symbols under a stated rule for the square roots. Only normalised series are canonical; the note normalises series by their constant terms, so only the $z$-dependence of the symbols matters, and it is always explicit. The convention $\re^{F/z}$ becomes the convention $\re^{-\cW/z}$, resp.\ $\re^{-\hat\cW/\zh}$, of \cref{sec:con} through $F=-\cW$, resp.\ $F=-\hat\cW$ with $z$ renamed $\zh$.
\end{defn}

% Formal stationary phase, part 2 of 5: universality, constants, homogeneity (end of sec:fsp-def); exact forms, Gaussian
% perturbation, products, formal implicit functions (sec:fsp-basic). Continued in part 3.

\begin{lem}[Universality of the coefficients]\genusstatus{AI-proposed, unchecked}
\label{lem:fsp-univ}
Let $c$ be a non-degenerate critical point of $F$ over $A$ and $a(c+t)=\sum_{m\geq m_0}a_m(t)z^m$, $a_m\in A\fps{t}$. Then
\begin{equation}
\label{eq:fsp-univ}
\coef{z^k}\,\SPn_c[F;a]=\sum_{m=m_0}^{k}\ \sum_{r=0}^{2(k-m)}\frac1{r!}\,\bigl\langle\,[F_{\geq3}^{\,r}\,a_m]_{2(k-m+r)}\bigr\rangle_C ,
\end{equation}
$[P]_d$ denoting the homogeneous part of degree $d$ in $t$. So $\coef{z^k}\SPn_c[F;a]$ is a polynomial over $\bbQ$ in the entries of $C$, the Taylor coefficients of $F$ at $c$ of orders $3,\dots,2(k-m_0)+2$ and those of the $a_m$ of orders $\leq2(k-m)$, linear in the latter. Consequently: (i) a ring homomorphism $\rho\colon A\to A'$, applied to the coefficients, maps $\SPn_c[F;a]$ to $\SPn_{\rho(c)}[\rho F;\rho a]$; (ii) if $A$ is a ring of holomorphic functions or of formal power series in some parameters, every coefficient is such a function, resp.\ series.
\end{lem}

\begin{proof}
The terms $F_{\geq3}^{\,r}a_mz^{m-r}/r!$ form a summable family with sum $\exp(F_{\geq3}/z)\,a(c+t)$, so $E_{zC}$ applies termwise. A monomial $t^\alpha$ of $F_{\geq3}^{\,r}a_m$ contributes to $z^k$ if and only if $|\alpha|=2(k-m+r)$; as $F_{\geq3}^{\,r}\in(t)^{3r}$, this forces $r\leq2(k-m)$ and $m\leq k$. A contributing term with factors of $F_{\geq3}$ of degrees $l_1,\dots,l_r\geq3$ and a monomial of $a_m$ of degree $j$ has $l_1+\dots+l_r+j=2(k-m+r)$, so $l_i\leq2(k-m)+2$ and $j\leq2(k-m)$. Moments are polynomials in the entries of $C$ (\cref{lem:fsp-ring}~(iii)). For (i), $\rho$ maps the Taylor data at $c$ to those at $\rho(c)$, and $C$ to $-\rho(\mathsf H)^{-1}$; (ii) follows.
\end{proof}

\begin{lem}[Constants]\genusstatus{AI-proposed, unchecked}
\label{lem:fsp-const}
For $\kappa\in A$, $\SPn_c[F+\kappa;a]=\SPn_c[F;a]$, and the full symbol of $\re^{(F+\kappa)/z}a\,\dd^nx$ is $[\re^{\kappa/z}]$ times that of $\re^{F/z}a\,\dd^nx$; so $\SPn_c[F;a]$ depends only on the Taylor coefficients of order $\geq2$ of $F$ at $c$ and on the Taylor series of $a$ at $c$.
\end{lem}

\begin{proof}
$\mathsf H$ and $F_{\geq3}$ do not change; the last assertion is \eqref{eq:fsp-univ}.
\end{proof}

\begin{lem}[Homogeneity]\genusstatus{AI-proposed, unchecked}
\label{lem:fsp-homog}
For $\mu\in A^\times$, $\SPn_c[\mu F;a(\cdot\,;z)](z)=\SPn_c[F;a(\cdot\,;\mu z')](z')$ at $z'=z/\mu$, i.e.\ after the isomorphism $A\lser{z}\to A\lser{z'}$, $z\mapsto\mu z'$; and, with the branch rule $\det(-\mu\mathsf H)^{1/2}:=\mu^{n/2}\det(-\mathsf H)^{1/2}$,
\[
[\re^{\mu F(c)/z}]\,(2\pi z)^{n/2}\,[\det(-\mu\mathsf H)^{-1/2}]=[\re^{F(c)/z'}]\,(2\pi z')^{n/2}\,[\det(-\mathsf H)^{-1/2}].
\]
\end{lem}

\begin{proof}
For $\mu F$ the Hessian is $\mu\mathsf H$, the covariance $C/\mu$ and the cubic part $\mu F_{\geq3}$. As $\langle t^\alpha\rangle_C$ is homogeneous of degree $|\alpha|/2$ in $C$, $E_{zC/\mu}[t^\alpha z^j]=\mu^{-|\alpha|/2}\langle t^\alpha\rangle_Cz^{j+|\alpha|/2}$, which at $z=\mu z'$ equals $\langle t^\alpha\rangle_C\mu^jz'^{\,j+|\alpha|/2}=E_{z'C}[t^\alpha(\mu z')^j]$. Both sides are linear, weight-preserving and compatible with summable sums, so $E_{zC/\mu}[P(t,z)]|_{z=\mu z'}=E_{z'C}[P(t,\mu z')]$ on $\cR_n(A)$; take $P=\exp(\mu F_{\geq3}/z)a(c+t;z)$, for which $P(t,\mu z')=\exp(F_{\geq3}/z')a(c+t;\mu z')$. The symbols agree term by term.
\end{proof}

%-------------------------------------------------------------------------------
\subsection{Exact forms, Gaussian perturbation, implicit functions and parameters}
\label{sec:fsp-basic}

\begin{lem}[Exact forms]\genusstatus{AI-proposed, unchecked}
\label{lem:fsp-ibp}
Let $c$ be a non-degenerate critical point of $F$ over $A$ and $h_1,\dots,h_n$ amplitudes with $\exp(F_{\geq3}/z)h_i\in\cR_n(A)$ (e.g.\ $h_i\in A\fps{x-c}\lser{z}$). Then $\SPn_c\bigl[F;\sum_i(\de_ih_i+z^{-1}(\de_iF)h_i)\bigr]=0$, i.e.\ $\mathrm{SP}_c[\dd(\re^{F/z}\iota)]=0$ for $\iota=\sum_i(-1)^{i-1}h_i\,\dd x_1\cdots\widehat{\dd x_i}\cdots\dd x_n$. The same holds for every sum of such expressions that is summable in $\cR_n(A)$ after multiplication by $\exp(F_{\geq3}/z)$.
\end{lem}

\begin{proof}
First, for every $g\in\cR_n(A)$ and every $i$,
\begin{equation}
\label{eq:fsp-gibp}
E_{zC}\bigl[\de_ig+z^{-1}(\mathsf Ht)_i\,g\bigr]=0\qquad\text{(Gaussian integration by parts).}
\end{equation}
Multiplying by $C$ and using $C\mathsf H=-1$, \eqref{eq:fsp-gibp} for all $i$ is equivalent to $E_{zC}[t_kg]=z\sum_jC_{kj}E_{zC}[\de_jg]$ for all $k$. For $g=t^\beta z^l$ both sides equal $z^{l+(|\beta|+1)/2}\sum_jC_{kj}\beta_j\langle t^{\beta-e_j}\rangle_C$ by the Wick recursion, and both sides are linear and compatible with summable sums. Now put $g_i:=\exp(F_{\geq3}/z)h_i$. Then $\de_ig_i=\exp(F_{\geq3}/z)(\de_ih_i+z^{-1}(\de_iF_{\geq3})h_i)$ and $\de_iF(c+t)=(\mathsf Ht)_i+\de_iF_{\geq3}(t)$, so $\exp(F_{\geq3}/z)\sum_i(\de_ih_i+z^{-1}(\de_iF)h_i)=\sum_i(\de_ig_i+z^{-1}(\mathsf Ht)_ig_i)$, whose image under $E_{zC}$ vanishes; summable sums are treated termwise.
\end{proof}

\begin{lem}[Gaussian perturbation]\genusstatus{AI-proposed, unchecked}
\label{lem:fsp-pg}
Let $R$ be a commutative $\bbQ$-algebra, complete and separated for an ideal $I$, $\mathsf H\in GL_n(R)$ symmetric, $C:=-\mathsf H^{-1}$, $\mathsf K$ symmetric with entries in $I$, and $C_{\mathsf K}:=-(\mathsf H+\mathsf K)^{-1}$.

(i) For $P\in R[t]$, $\langle\exp(\frac12t^{\mathrm T}\mathsf Kt)P\rangle_C=\det(1+\mathsf H^{-1}\mathsf K)^{-1/2}\langle P\rangle_{C_{\mathsf K}}$, both sides converging $I$-adically, the power being the binomial series.

(ii) (Infinitesimal form.) For any symmetric $\mathsf K'$ over $R$, a variable $s$, $C_s:=-(\mathsf H+s\mathsf K')^{-1}$ and $d_s:=\det(1+s\mathsf H^{-1}\mathsf K')^{-1/2}$ in $R\fps{s}$: $\frac{\dd}{\dd s}(d_s\langle P\rangle_{C_s})=d_s\langle\frac12t^{\mathrm T}\mathsf K't\,P\rangle_{C_s}$.

(iii) (With $z$ and outer variables.) Let $\mathsf H\in GL_n(A)$ be symmetric, $C=-\mathsf H^{-1}$, $y=(y_1,\dots,y_m)$ outer variables, $\mathsf K$ symmetric with entries in the ideal $(y)$ of $A\fps{y}$, and $\mathsf H':=\mathsf H+\mathsf K$. Then every term of $t^{\mathrm T}\mathsf Kt/(2z)$ has weight $\geq1/2$, and for $P\in\cR_{n+m}(A)$
\[
E^t_{zC}\bigl[\exp\bigl(t^{\mathrm T}\mathsf Kt/(2z)\bigr)P\bigr]=\bigl(\det\mathsf H'/\det\mathsf H\bigr)^{-1/2}E^t_{-z\mathsf H'^{-1}}[P]\qquad\text{in }\cR_m(A).
\]
The same holds over $A_0\fps{q}$ for $\mathsf K$ with entries in $(q)$ and $P\in\cR_n(A_0)\fps{q}$, coefficientwise in $q$ (in the sense of \cref{def:fsp-ring}).
\end{lem}

\begin{proof}
(i), (ii). By the generating function $\langle P\rangle_C=[\exp(\frac12\sum_{ij}C_{ij}\de_i\de_j)P](0)$ (on $t^\alpha$ the operator removes the variables in pairs). Let $Q:=\frac12t^{\mathrm T}\mathsf K't$ and $\rho_s(P):=d_s\langle P\rangle_{C_s}$. As $\dd C_s/\dd s=C_s\mathsf K'C_s$, $\frac{\dd}{\dd s}\langle P\rangle_{C_s}=\frac12\sum_{ij}(C_s\mathsf K'C_s)_{ij}\langle\de_i\de_jP\rangle_{C_s}$; applying the Wick recursion $\langle t_iP\rangle_C=\sum_lC_{il}\langle\de_lP\rangle_C$ twice,
\[
\langle QP\rangle_{C_s}=\tfrac12\tr(\mathsf K'C_s)\,\langle P\rangle_{C_s}+\tfrac12\textstyle\sum_{lm}(C_s\mathsf K'C_s)_{lm}\langle\de_l\de_mP\rangle_{C_s};
\]
and $\frac{\dd}{\dd s}\log d_s=-\frac12\tr((\mathsf H+s\mathsf K')^{-1}\mathsf K')=\frac12\tr(C_s\mathsf K')$. Together $\frac{\dd}{\dd s}\rho_s(P)=\rho_s(QP)$, which is (ii). For (i) take $\mathsf K'=\mathsf K$ and $\Lambda_s(P):=\langle\exp(sQ)P\rangle_C=\sum_ks^k\langle Q^kP\rangle_C/k!$; then $\frac{\dd}{\dd s}\Lambda_s(P)=\Lambda_s(QP)$ and $\Lambda_0=\rho_0$, so all $s$-derivatives at $0$ agree ($\langle Q^kP\rangle_C$) and $\Lambda_s(P)=\rho_s(P)$ in $R\fps{s}$. The $s^k$-coefficients of both sides, of $d_s$ and of $C_s=-\sum_r(-s)^r(\mathsf H^{-1}\mathsf K)^r\mathsf H^{-1}$ lie in $I^k$, and $s=1$ is a ring homomorphism from such series to the complete ring $R$; it gives (i).
(iii) The terms of $t^{\mathrm T}\mathsf Kt/(2z)$ are multiples of $t_it_jy^\gamma z^{-1}$, $|\gamma|\geq1$. For $P=t^\alpha y^\beta z^j$, the term $(t^{\mathrm T}\mathsf Kt/(2z))^kt^\alpha/k!$ has degree $2k+|\alpha|$ in $t$, so $E^t_{zC}[\exp(t^{\mathrm T}\mathsf Kt/(2z))t^\alpha y^\beta z^j]=\langle\exp(\frac12t^{\mathrm T}\mathsf Kt)t^\alpha\rangle_Cy^\beta z^{j+|\alpha|/2}$, which by (i) over $R=A\fps{y}$, $I=(y)$ (with $\det(1+\mathsf H^{-1}\mathsf K)=\det\mathsf H'/\det\mathsf H$, $C_{\mathsf K}=-\mathsf H'^{-1}$) is the right side for this $P$. Both sides are linear and compatible with summable sums (the exponential has terms of weight $\geq0$ only), so the identity extends to $\cR_{n+m}(A)$. Over $A_0\fps q$ argue modulo $q^N$: over $B_N=A_0[q]/(q^N)$ the ideal $(q)$ is nilpotent, (i) applies with $R=B_N$ and $I=(q)$, and the monomial argument above gives the identity on $\cR_n(B_N)$.
\end{proof}

\begin{lem}[Products]\genusstatus{AI-proposed, unchecked}
\label{lem:fsp-prod}
Split the variables as $(t,y)$ and let $C=\diag(C_t,C_y)$ with constant blocks. Then, for $P\in\cR_{n+m}(A)$, $E_{zC}[P]=E^y_{zC_y}\bigl[E^t_{zC_t}[P]\bigr]$. Hence, for phases and amplitudes in separate variables,
\[
\SPn_{(c_1,c_2)}\bigl[F_1(t)+F_2(y);\,a_1(t)\,a_2(y)\bigr]=\SPn_{c_1}[F_1;a_1]\,\SPn_{c_2}[F_2;a_2],
\]
and the full symbols multiply, with $\det(-\mathsf H)^{1/2}:=\det(-\mathsf H_1)^{1/2}\det(-\mathsf H_2)^{1/2}$.
\end{lem}

\begin{proof}
For block-diagonal $C$ the generating function factorises, so $\langle t^\alpha y^\beta\rangle_C=\langle t^\alpha\rangle_{C_t}\langle y^\beta\rangle_{C_y}$; both sides are linear and compatible with summable sums. For separated data the integrand is the product of $\exp(F_{1,\geq3}/z)a_1$, a function of $t$, and $\exp(F_{2,\geq3}/z)a_2$, a function of $y$; and $E^t[P_1(t)P_2(y)]=E^t[P_1]\,P_2$ with $E^t[P_1]\in A\lser{z}$.
\end{proof}

\begin{lem}[Formal implicit function theorem]\genusstatus{AI-proposed, unchecked}
\label{lem:fsp-ift}
Let $B$ be a commutative $\bbQ$-algebra, complete and separated for an ideal $I$ (e.g.\ $B=A\fps{y}$, $I=(y)$, or $B=A_0\fps{q}$, $I=(q)$), and $F\in B\fps{t}$ with $\nabla F(0)\in I^n$ and $\mathsf H_0:=\Hess F(0)\in GL_n(B)$. Then there is a unique $c\in I^n$ with $\nabla F(c)=0$, and $\Hess F(c)$ is invertible.
\end{lem}

\begin{proof}
Substitution of elements of $I$ into power series converges, $B$ being complete. For $c\in I^n$ put $\Psi(c):=c-\mathsf H_0^{-1}\nabla F(c)$; it lies in $I^n$, because $\nabla F(c)=\nabla F(0)+\mathsf H_0c+O(c^2)$ and $\nabla F(0)\in I^n$. Writing $\nabla F(c)-\nabla F(c')=\mathsf N(c,c')(c-c')$ with difference quotients, $\mathsf N(c,c')\equiv\mathsf H_0$ modulo $I$, one gets $\Psi(c)-\Psi(c')=-\mathsf H_0^{-1}(\mathsf N(c,c')-\mathsf H_0)(c-c')$, a matrix with entries in $I$ times $c-c'$; so $c-c'\in I^r$ implies $\Psi(c)-\Psi(c')\in I^{r+1}$. The iterates $c_{k+1}:=\Psi(c_k)$, $c_0=0$, satisfy $c_{k+1}-c_k\in I^{k+1}$, so they converge to some $c\in I^n$ with $\Psi(c)=c$, i.e.\ $\nabla F(c)=0$. Two such zeros $c,c'$ satisfy $c-c'=\Psi(c)-\Psi(c')$, hence $c-c'\in\bigcap_rI^r=0$. Finally $\Hess F(c)\equiv\mathsf H_0$ modulo $I$, and a matrix over $B$ invertible modulo $I$ is invertible.
\end{proof}

% Formal stationary phase, part 3 of 5 (sec:fsp-basic, continued): translations, moving the base point, parameters,
% derivatives in a parameter, the stationary phase at a generic base point. Continued in part 4.

\begin{lem}[Translations]\genusstatus{AI-proposed, unchecked}
\label{lem:fsp-transl}
Let $c$ be a non-degenerate critical point of $F$ over $A$, $a$ an amplitude and $p\in A^n$ small in one of two senses: (a) $A=A_0\fps{q}$, $F$ and $c$ are defined over $A_0$, $a\in A_0\fps{x-c}\lser{z}\fps{q}$ and $p\in(qA)^n$; or (b) $p\in(zA\fps{z})^n$ and $a\in A\fps{x-c}\lser{z}$. Then the translate $b_p(t):=\exp\bigl((F(c+t+p)-F(c+t))/z\bigr)a(c+t+p)$ (the formal analogue of $x\mapsto x+p$ applied to $\re^{F/z}a$) is an amplitude, in $A_0\fps{t}\lser{z}\fps{q}$, resp.\ $A\fps{t}\lser{z}$, and $\SPn_c[F;b_p]=\SPn_c[F;a]$.
\end{lem}

\begin{proof}
In case (a), $F(c+t+p)-F(c+t)$ has coefficients in $qA$, so the exponential converges $q$-adically, each $q$-coefficient lying in $A_0\fps{t}[z^{-1}]$. In case (b), with $p=zv$, $(F(c+t+zv)-F(c+t))/z=\nabla F(c+t)\cdot v+O(z)$ lies in the ideal $(t,z)$ of $A\fps{t,z}$ because $\nabla F(c+t)=O(t)$, so the exponential converges in $A\fps{t,z}$. Let $s$ be a variable and $B(s):=b_{sp}$. At fixed degree in $t$ and powers of $z$ (and $q$), the coefficient of $B(s)$ is a polynomial in $s$: in case (a) each $p_i$ carries a factor $q$; in case (b) each factor $\nabla F(c+t)\cdot sv$ of the exponential has degree $\geq1$ in $t$, while every other factor, and every term of the Taylor expansion of $a(c+t+szv)$ that involves $s$, carries a positive power of $z$. By \cref{lem:fsp-univ} the coefficients of $\SPn_c[F;B(s)]$ are polynomials in $s$, with derivative $\SPn_c[F;\dd B/\dd s]$. Differentiating $B(s)$ and comparing with $\de_iB(s)$, $\de_i=\de/\de t_i$, gives $\dd B/\dd s=\sum_ip_i\bigl(\de_iB+z^{-1}(\de_iF)(c+t)B\bigr)$; by \cref{lem:fsp-ibp} with $h_i:=p_iB(s)$ (coefficientwise in $q$ in case (a)) the derivative vanishes. A polynomial over a $\bbQ$-algebra with zero derivative is constant: $\SPn_c[F;B(1)]=\SPn_c[F;B(0)]$.
\end{proof}

\begin{cor}[Moving the base point]\genusstatus{AI-proposed, unchecked}
\label{cor:fsp-movebase}
Let $A:=A_0\fps{q}$, let $F_0\in A_0\fps{x-c_0}$ have the non-degenerate critical point $c_0$ over $A_0$, with Hessian $\mathsf H_0$, let $F:=F_0+F_1$ with $F_1\in qA_0\fps{x-c_0}\fps{q}$, and $a\in A_0\fps{x-c_0}\lser{z}\fps{q}$. Then $F$ has a unique critical point $c\in c_0+(qA)^n$, non-degenerate with Hessian $\mathsf H_c$. For amplitudes $b$ put
\[
\mathcal S_q[F;b]:=\exp\bigl((F(c)-F_0(c_0))/z\bigr)\,(\det\mathsf H_c/\det\mathsf H_0)^{-1/2}\,\SPn_c[F;b]
\]
(the power is a binomial series, as $\det\mathsf H_c/\det\mathsf H_0\in1+qA$, and the exponential converges, as $(F(c)-F_0(c_0))/z$ lies in $qA_0\fps{q}\,z^{-1}$). Then $\SPn_{c_0}[F_0;\re^{F_1/z}a]=\mathcal S_q[F;a]$ in $A_0\lser{z}\fps{q}$, the left side being computed coefficientwise in $q$. With the branch rule $\det(-\mathsf H_c)^{1/2}:=\det(-\mathsf H_0)^{1/2}(\det\mathsf H_c/\det\mathsf H_0)^{1/2}$, $\mathcal S_q[F;a]$ is $\mathrm{SP}_c[\re^{F/z}a\,\dd^nx]$ divided by the $q$-independent symbols $[\re^{F_0(c_0)/z}](2\pi z)^{n/2}[\det(-\mathsf H_0)^{-1/2}]$.
\end{cor}

\begin{proof}
\Cref{lem:fsp-ift} over $B=A$, $I=(q)$, gives $c$: $\nabla F(c_0)=\nabla F_1(c_0)\in(q)^n$, and $\Hess F(c_0)\in\mathsf H_0+q\,\mathfrak{gl}_n(A)$ is invertible. Put $p:=c-c_0$. By \cref{lem:fsp-transl}~(a) for $F_0$ and the amplitude $\re^{F_1/z}a$, the left side is $\SPn_{c_0}[F_0;b_p]$, $b_p(t)=\exp((F(c+t)-F_0(c_0+t))/z)\,a(c+t)$. With $F(c+t)=F(c)+\frac12t^{\mathrm T}\mathsf H_ct+F_{c,\geq3}(t)$ ($c$ is critical) and $F_0(c_0+t)=F_0(c_0)+\frac12t^{\mathrm T}\mathsf H_0t+F_{0,\geq3}(t)$, in $\cR_n(A_0)\fps{q}$, that is, modulo every $q^N$ (\cref{def:fsp-ring}; both sides are the exponential of $(F(c+t)-F_0(c_0)-\frac12t^{\mathrm T}\mathsf H_0t)/z$),
\[
\re^{F_{0,\geq3}/z}\,\re^{(F(c+t)-F_0(c_0+t))/z}=\re^{(F(c)-F_0(c_0))/z}\,\re^{t^{\mathrm T}(\mathsf H_c-\mathsf H_0)t/(2z)}\,\re^{F_{c,\geq3}/z}.
\]
The scalar factor comes out of $E_{-z\mathsf H_0^{-1}}$. As $\mathsf H_c-\mathsf H_0$ has entries in $qA$, \cref{lem:fsp-pg}~(iii) gives, for $P\in\cR_n(A_0)\fps{q}$,
\[
E_{-z\mathsf H_0^{-1}}\bigl[\re^{t^{\mathrm T}(\mathsf H_c-\mathsf H_0)t/(2z)}P\bigr]=(\det\mathsf H_c/\det\mathsf H_0)^{-1/2}\,E_{-z\mathsf H_c^{-1}}[P],
\]
and $E_{-z\mathsf H_c^{-1}}[\re^{F_{c,\geq3}/z}a(c+t)]=\SPn_c[F;a]$.
\end{proof}

\begin{lem}[Parameters]\genusstatus{AI-proposed, unchecked}
\label{lem:fsp-param}
(i) (Dependence on parameters.) If the Taylor coefficients at $c$ of $F$ (orders $\geq2$) and of $a$ lie in a $\bbQ$-subalgebra $A'\subset A$ and $\det\mathsf H\in A'^\times$, all coefficients of $\SPn_c[F;a]$ lie in $A'$; rational data give rational coefficients, with poles only where the data have poles or $\det\mathsf H=0$. Over $A=A_0\fps{q}$, $q=(q_1,\dots,q_p)$, if $F|_{q=0}$ has the non-degenerate critical point $c_0$, then $F$ has a unique critical point $c(q)\in c_0+(q)A^n$, non-degenerate, and the coefficients of $\SPn_{c(q)}[F;a]$, $F(c(q))$ and $\det\Hess F(c(q))$ lie in $A_0\fps{q}$. If $F$, $a$ are holomorphic near $(c_0,q_0)$ and $c(q)$ is a holomorphic family of non-degenerate critical points with $c(q_0)=c_0$, these are holomorphic near $q_0$.

(ii) (Taylor expansion and evaluation.) In the holomorphic case, the Taylor series at $q_0$ of each coefficient is the corresponding coefficient over $\bbC\fps{q-q_0}$ for the Taylor-expanded data, at the critical point of the formal case, and its value at $q_0$ is that of $\SPn_{c_0}[F(\cdot\,;q_0);a(\cdot\,;q_0)]$. In the formal case, $q\mapsto0$ maps $\SPn_{c(q)}[F;a]$ to $\SPn_{c_0}[F|_{q=0};a|_{q=0}]$.

(iii) (Perturbation in a parameter.) \Cref{cor:fsp-movebase}: the $q$-expansion of the stationary phase at the moving critical point, with its two prefactors, is that at $c_0$ with $\re^{F_1/z}$ in the amplitude; for $F_1=qg$ and $a$ free of $q$, $\coef{q^d}\SPn_{c_0}[F_0;\re^{qg/z}a]=\SPn_{c_0}[F_0;(g/z)^da]/d!$.

(iv) (A $z$-dependent part of the phase.) Let $c$ be a non-degenerate critical point of $F$ over $A$, $a\in A\fps{x-c}\lser z$ an amplitude (free of $\varepsilon$), $g\in A\fps{x-c}$, $g(c)=0$, $\varepsilon$ a variable and $F_\varepsilon:=F+\varepsilon g$, with its unique critical point $c(\varepsilon)\in c+(\varepsilon A\fps{\varepsilon})^n$. Then $\mathcal S_\varepsilon[F_\varepsilon;a]=\SPn_c[F;\re^{\varepsilon g/z}a]$; giving $\varepsilon$ the weight $1$, this series has finitely many non-zero terms $\varepsilon^dz^k$ of each bounded weight $d+k$, the substitution $\varepsilon\mapsto z$ is well defined on such series, and $[\mathcal S_\varepsilon[F_\varepsilon;a]]_{\varepsilon=z}=\SPn_c[F;\re^ga]$. So the stationary phase with the $z$-dependent phase $F+zg$ is defined and equals the one with $zg$ moved into the amplitude as $\re^g$.
\end{lem}

\begin{proof}
(i) The first assertion is \eqref{eq:fsp-univ}, as $C=-\operatorname{adj}(\mathsf H)/\det\mathsf H$. In the formal case \cref{lem:fsp-ift} over $A$, $I=(q)$, gives $c(q)$, and the data at $c(q)$ lie in $A$. In the holomorphic case $\de_x^\alpha F(c(q);q)$, $\de_x^\alpha a(c(q);q)$ and the covariance are holomorphic in $q$. (ii) The Taylor map and the evaluations are ring homomorphisms (\cref{lem:fsp-univ}~(i)); the Taylor series of $c(q)$ is a critical point of the expanded phase, hence the formal one by uniqueness (\cref{lem:fsp-ift}). (iii) $\re^{qg/z}a=\sum_dq^d(g/z)^da/d!$. (iv) \Cref{lem:fsp-ift} over $A\fps{\varepsilon}$ gives $c(\varepsilon)$, as $\nabla F_\varepsilon(c)=\varepsilon\nabla g(c)$; the first identity is \cref{cor:fsp-movebase} with $q=\varepsilon$, $F_0=F$, $F_1=\varepsilon g$. In $\cR_{n;\varepsilon}(A)$ every term of $\varepsilon g/z$ and of $F_{\geq3}/z$ has weight $\geq1/2$ ($g=O(t)$), so $\exp(F_{\geq3}/z)\exp(\varepsilon g/z)a(c+t)$ lies there; it is the same formal sum as in $\cR_n(A)\fps{\varepsilon}$, $E_{zC}$ preserves the weight, and each coefficient of $\varepsilon^dz^k$ is the same finite sum in both readings, which gives the finiteness. The substitution $t^\alpha\varepsilon^dz^j\mapsto t^\alpha z^{d+j}$ maps $\cR_{n;\varepsilon}(A)$ to $\cR_n(A)$ (each output monomial comes from finitely many inputs of its weight), is a ring homomorphism compatible with summable sums, maps $\exp(\varepsilon g/z)$ to $\re^g$, and commutes with $E_{zC}$ (both composites send $t^\alpha\varepsilon^dz^j$ to $\langle t^\alpha\rangle_Cz^{d+j+|\alpha|/2}$).
\end{proof}

\begin{lem}[Derivatives in a parameter]\genusstatus{AI-proposed, unchecked}
\label{lem:fsp-dq}
In the situation of \cref{cor:fsp-movebase}, for every amplitude $a\in A_0\fps{x-c_0}\lser{z}\fps{q}$,
\[
\frac{\dd}{\dd q}\,\mathcal S_q[F;a]=\mathcal S_q\bigl[F;\,z^{-1}(\de_qF)\,a+\de_qa\bigr]\qquad\text{in }A_0\lser{z}\fps{q},
\]
i.e.\ $\frac{\dd}{\dd q}\mathrm{SP}[\re^{F/z}a\,\dd^nx]=\mathrm{SP}[\de_q(\re^{F/z}a)\,\dd^nx]$; the same holds for $\mathrm D:=(q_0+q)\,\dd/\dd q$, $q_0\in A_0$, in place of $\de_q$. If $\mathrm DF=\varphi$ does not depend on $q$ and $\mathrm Da=0$, then $P(z\mathrm D)\,\mathcal S_q[F;a]=\mathcal S_q[F;P(\varphi)a]$ for every polynomial $P$ over $A_0[z]$.
\end{lem}

\begin{proof}
By \cref{cor:fsp-movebase}, $\mathcal S_q[F;a]=\SPn_{c_0}[F_0;\re^{F_1/z}a]$ with a $q$-independent functional, applied coefficientwise in $q$; so $\dd/\dd q$ acts on the amplitude, $\frac{\dd}{\dd q}(\re^{F_1/z}a)=\re^{F_1/z}(z^{-1}(\de_qF)a+\de_qa)$, and \cref{cor:fsp-movebase} for this amplitude gives the claim. For $\mathrm D$ multiply by $q_0+q$. If $\mathrm DF=\varphi$ is $q$-independent and $\mathrm Da=0$, then $\mathrm D(\varphi^ka)=0$ and $z\mathrm D\,\mathcal S_q[F;\varphi^ka]=\mathcal S_q[F;\varphi^{k+1}a]$; induction and linearity over $A_0[z]$ give the clause for $P$ (without $\mathrm Da=0$, terms $z\mathrm D(\varphi^ka)$ appear).
\end{proof}

\begin{lem}[The stationary phase at a generic base point]\genusstatus{AI-proposed, unchecked}
\label{lem:fsp-generic}
Let $\delta(\sigma):=(\sigma^6-\sigma)/5$, so that $\AL=\bbQ[L,\delta(L)^{-1}]$ and $\Dq L=\delta(L)$, and let $\bbA_\sigma:=\bbQ[\sigma,\delta(\sigma)^{-1}]$ be the same ring in a variable $\sigma$; for $f\in\bbA_\sigma$ let $f(L+t)\in\AL\fps{t}$ be its Taylor expansion at $\sigma=L$. Let $S$ be a phase at $\sigma=L$ over $\AL$ with $S'(L)=0$ and $S''(L+t)=-1/\delta(L+t)$, $\varepsilon$ a variable and $F(\sigma;\varepsilon):=S(\sigma)+\varepsilon\sigma$. Let $\mathsf r(\sigma):=5\sigma/(\sigma^5-1)\in\bbA_\sigma^\times$, and for $\sigma_\ast\in\{L,\sigma_\varepsilon\}$ let $(\mathsf r(\sigma)/\mathsf r(\sigma_\ast))^{1/2}$ be the square root with constant term $1$ of the expansion of $\mathsf r(\sigma)/\mathsf r(\sigma_\ast)$ at $\sigma_\ast$.

(i) $F(\cdot\,;\varepsilon)$ has a unique critical point $\sigma_\varepsilon\in L+\varepsilon\AL\fps{\varepsilon}$, non-degenerate, and $\sigma_\varepsilon=\exp(\varepsilon\Dq)(L)$, where $\exp(\varepsilon\Dq):=\sum_m\varepsilon^m\Dq^m/m!$; the ring homomorphism $\AL\to\AL\fps{\varepsilon}$, $L\mapsto\sigma_\varepsilon$, is $\exp(\varepsilon\Dq)$.

(ii) For $b\in\bbA_\sigma\fps{z}$, with $\exp(\varepsilon\Dq)$ acting coefficientwise in $z$,
\[
\SPn_{\sigma_\varepsilon}\bigl[F(\cdot\,;\varepsilon);(\mathsf r(\sigma)/\mathsf r(\sigma_\varepsilon))^{1/2}b\bigr]=\exp(\varepsilon\Dq)\,\SPn_L\bigl[S;(\mathsf r(\sigma)/\mathsf r(L))^{1/2}b\bigr].
\]

(iii) For $a:=(\mathsf r(\sigma)/\mathsf r(L))^{1/2}b$, $b\in\bbA_\sigma\fps{z}$, and $\Dq$ acting coefficientwise in $z$,
\begin{equation}
\label{eq:fsp-generic}
(z\Dq+L)\bigl(L\cdot\SPn_L[S;a]\bigr)=L\cdot\SPn_L[S;\sigma a]\qquad\text{in }\AL\lser{z},
\end{equation}
where $\sigma a$ is the amplitude $(L+t)\,a(L+t,z)$.
\end{lem}

\begin{proof}
(i) $F(L+t;\varepsilon)=S(L)+\varepsilon L+\varepsilon t+\sum_{k\geq2}S^{(k)}(L)t^k/k!$ has gradient $\varepsilon$ and Hessian $-1/\delta(L)\in\AL^\times$ at $t=0$, so \cref{lem:fsp-ift} over $\AL\fps{\varepsilon}$ gives $\sigma_\varepsilon=L+t_\varepsilon$, $t_\varepsilon\in\varepsilon\AL\fps{\varepsilon}$. For $f\in\bbA_\sigma$, the substitution $t=t_\varepsilon$ in $f(L+t)$ is the ring homomorphism $\bbA_\sigma\to\AL\fps{\varepsilon}$, $\sigma\mapsto\sigma_\varepsilon$ (well defined, $\delta(\sigma_\varepsilon)\equiv\delta(L)$ modulo $\varepsilon$ being a unit): both are ring homomorphisms agreeing on $\sigma$. So $S''(\sigma_\varepsilon)=-1/\delta(\sigma_\varepsilon)$, and differentiating $S'(\sigma_\varepsilon)+\varepsilon=0$ gives $\sigma_\varepsilon'=\delta(\sigma_\varepsilon)$, $\sigma_0=L$. On the other hand $\exp(\varepsilon\Dq)$ is a ring homomorphism (Leibniz rule for $\Dq^m$) with $\frac{\dd}{\dd\varepsilon}\exp(\varepsilon\Dq)(f)=\exp(\varepsilon\Dq)(\Dq f)$, so $\tilde\sigma_\varepsilon:=\exp(\varepsilon\Dq)(L)$ satisfies $\tilde\sigma_\varepsilon'=\delta(\tilde\sigma_\varepsilon)$, $\tilde\sigma_0=L$. A solution $y\in\AL\fps{\varepsilon}$ of $y'=\delta(y)$, $y(0)=L$, is unique, since $(m+1)y_{m+1}=\coef{\varepsilon^m}\delta(y)$ involves only $y_0,\dots,y_m$. So $\sigma_\varepsilon=\tilde\sigma_\varepsilon$, and the two homomorphisms agree on $L$, hence on $\AL$.
(ii) For $f\in\bbA_\sigma$, the expansion $f(\sigma_\varepsilon+t)$ is $\exp(\varepsilon\Dq)$ applied coefficientwise to $f(L+t)$: both are ring homomorphisms $\bbA_\sigma\to\AL\fps{\varepsilon}\fps{t}$ sending $\sigma$ to $\sigma_\varepsilon+t$, by (i). This applies to the derivatives of $-1/\delta$ (the Taylor coefficients of order $\geq2$ of $F(\cdot\,;\varepsilon)$ at $\sigma_\varepsilon$), to the covariance $\delta(\sigma_\varepsilon)$, to the coefficients of $b$ and to $\mathsf r$; the binomial series has rational coefficients. \Cref{lem:fsp-univ}~(i) and \cref{lem:fsp-const} give (ii).
(iii) Put $N:=\SPn_L[S;a]$ and $N^\sigma:=\SPn_L[S;\sigma a]$. The amplitude $a$ is free of $\varepsilon$, and at $\sigma_\varepsilon$ it equals $(\mathsf r(\sigma_\varepsilon)/\mathsf r(L))^{1/2}(\mathsf r(\sigma)/\mathsf r(\sigma_\varepsilon))^{1/2}b$, the first factor being the constant $(\exp(\varepsilon\Dq)(\mathsf r(L))/\mathsf r(L))^{1/2}$. By linearity and (ii), and since the Hessian at $\sigma_\varepsilon$ is $\exp(\varepsilon\Dq)(-1/\delta(L))$, the series of \cref{cor:fsp-movebase} (base ring $\AL$, parameter $\varepsilon$, $F_0=S$, $F_1=\varepsilon\sigma$) is
\[
\mathcal S_\varepsilon[F;a]=\re^{(F(\sigma_\varepsilon;\varepsilon)-S(L))/z}\Bigl(\frac{\exp(\varepsilon\Dq)(\delta(L)\mathsf r(L))}{\delta(L)\mathsf r(L)}\Bigr)^{1/2}\exp(\varepsilon\Dq)(N)=\re^{(F(\sigma_\varepsilon;\varepsilon)-S(L))/z}\,\frac{\sigma_\varepsilon}L\,\exp(\varepsilon\Dq)(N),
\]
because $\delta(L)\mathsf r(L)=L^2$ and two elements of $1+\varepsilon\AL\fps{\varepsilon}$ with equal squares are equal. By \cref{lem:fsp-dq} with $\de_\varepsilon F=\sigma$, $\de_\varepsilon a=0$, the derivative at $\varepsilon=0$ of the left side is $z^{-1}N^\sigma$. On the right, $\frac{\dd}{\dd\varepsilon}F(\sigma_\varepsilon;\varepsilon)=\sigma_\varepsilon$ by criticality, so the three factors contribute $L/z$, $\Dq L/L$ and $\Dq N$. Hence $(L/z)N+(\Dq L/L)N+\Dq N=z^{-1}N^\sigma$; multiply by $zL$.
\end{proof}

% Formal stationary phase, part 4 of 5 (sec:fsp-change): the Piola identity, changes of variables, the uniformising map,
% iterated stationary phase. Continued in part 5.

%-------------------------------------------------------------------------------
\subsection{Changes of variables and iterated stationary phase}
\label{sec:fsp-change}

\begin{lem}[Piola identity]\genusstatus{AI-proposed, unchecked}
\label{lem:fsp-piola}
Let $B$ be a commutative $\bbQ$-algebra, $\psi\in B\fps{y}^n$, $y=(y_1,\dots,y_n)$, and $J:=D\psi=(\de\psi_i/\de y_k)_{ik}$. Then $\sum_k\de_k\operatorname{adj}(J)_{ki}=0$ for every $i$, where $\operatorname{adj}(J)_{ki}=(-1)^{i+k}\det J^{(i,k)}$ and $J^{(i,k)}$ is $J$ without row $i$ and column $k$.
\end{lem}

\begin{proof}
In the formal de Rham complex of $B\fps{y}$ the $(n-1)$-form $\omega_i:=\dd\psi_1\wedge\dots\wedge\widehat{\dd\psi_i}\wedge\dots\wedge\dd\psi_n$ is closed, each $\dd\psi_j$ being closed. Its coefficient on $\dd y_1\wedge\dots\widehat{\dd y_k}\dots\wedge\dd y_n$ is $\det J^{(i,k)}$, and $\dd(f\,\dd y_1\wedge\dots\widehat{\dd y_k}\dots\wedge\dd y_n)=(-1)^{k-1}\de_kf\,\dd y_1\wedge\dots\wedge\dd y_n$. So $0=\dd\omega_i=\sum_k(-1)^{k-1}\de_k\det J^{(i,k)}\,\dd y_1\wedge\dots\wedge\dd y_n$; multiply by $(-1)^{i+1}$.
\end{proof}

\begin{lem}[Change of variables]\genusstatus{AI-proposed, unchecked}
\label{lem:fsp-cov}
Let $F\in A\fps{x-c}$ have the non-degenerate critical point $c$ over $A$, with Hessian $\mathsf H$, and let $\phi\in c+\bigl((y-c')A\fps{y-c'}\bigr)^n$ be a formal map with $\phi(c')=c$ and $J:=D\phi(c')\in GL_n(A)$. Then $c'$ is a non-degenerate critical point of $F\circ\phi$, with Hessian $J^{\mathrm T}\mathsf HJ$, and for every amplitude $a\in A\fps{x-c}\lser z$
\[
\SPn_c[F;a]=\SPn_{c'}\bigl[F\circ\phi;\ (a\circ\phi)\det D\phi/\det J\bigr];
\]
for full symbols, $\mathrm{SP}_c[\re^{F/z}a\,\dd^nx]=\mathrm{SP}_{c'}[\re^{(F\circ\phi)/z}(a\circ\phi)\det(D\phi)\,\dd^ny]$ with the branch rule $\det(-J^{\mathrm T}\mathsf HJ)^{1/2}:=\det(J)\det(-\mathsf H)^{1/2}$. As $A$ is arbitrary, this covers analytic maps (through their Taylor series), formal maps and maps depending on parameters ($A=A_0\fps{u}$, or holomorphic in $u$ by \cref{lem:fsp-param}), including critical points that move with the parameter.
\end{lem}

\begin{proof}
$\nabla(F\circ\phi)(c')=J^{\mathrm T}\nabla F(c)=0$ and $\Hess(F\circ\phi)(c')=J^{\mathrm T}\mathsf HJ+\sum_i(\de_iF)(c)\Hess\phi_i(c')=J^{\mathrm T}\mathsf HJ$. Translating, let $c=c'=0$.
\emph{Linear maps.} For $\phi=J$ the covariance of $F\circ J$ is $J^{-1}CJ^{-\mathrm T}$ and its cubic part $F_{\geq3}\circ J$; \cref{lem:fsp-ring}~(v) gives $E_{zJ^{-1}CJ^{-\mathrm T}}[\re^{F_{\geq3}(Jy)/z}a(Jy)]=E_{zC}[\re^{F_{\geq3}/z}a]$, and $\det J/\det J=1$.
\emph{Reduction.} Write $\phi=J\circ\psi$, $\psi:=J^{-1}\phi$ (local letter), so $D\psi(0)=1$ and $\det D\phi=\det J\det D\psi$; with $\tilde F:=F\circ J$ and $\tilde a:=a\circ J$ it suffices to show $\SPn_0[\tilde F;\tilde a]=\SPn_0[\tilde F\circ\psi;(\tilde a\circ\psi)\det D\psi]$.
\emph{Maps tangent to the identity.} Write $\psi(y)=y+\chi(y)$, $\chi\in(y)^2$; for a variable $s$ put $\psi_s:=y+s\chi$, $J_s:=D\psi_s=1+sD\chi$ (invertible over $A[s]\fps{y}$), $b_s:=(\tilde a\circ\psi_s)\det J_s$ and $\cN(s):=\SPn_0[\tilde F\circ\psi_s;b_s]$. The phase $\tilde F\circ\psi_s$ has the critical point $0$ with the Hessian $\tilde{\mathsf H}$ of $\tilde F$ for all $s$, so $\cN(s)=E_{z\tilde C}[\exp(P_s/z)b_s]$ with $P_s:=\tilde F\circ\psi_s-\tilde F(0)-\frac12y^{\mathrm T}\tilde{\mathsf H}y$ and the fixed covariance $\tilde C=-\tilde{\mathsf H}^{-1}$. By \cref{lem:fsp-univ} each coefficient of $\cN$ is a polynomial in $s$, and $\cN'(s)=E_{z\tilde C}[\exp(P_s/z)(z^{-1}\de_s(\tilde F\circ\psi_s)b_s+\de_sb_s)]$. Put $\xi:=J_s^{-1}\chi$ (local letter). Then (a) $\de_s(\tilde F\circ\psi_s)=\sum_k\de_k(\tilde F\circ\psi_s)\xi_k$, both sides being $\sum_i(\de_i\tilde F)(\psi_s)\chi_i$, since $\de_k(\tilde F\circ\psi_s)=\sum_i(\de_i\tilde F)(\psi_s)(J_s)_{ik}$ and $J_s\xi=\chi$; and (b) $\de_sb_s=\sum_k\de_k(b_s\xi_k)$: with $\operatorname{adj}(J_s)=\det(J_s)J_s^{-1}$, $b_s\xi=(\tilde a\circ\psi_s)\operatorname{adj}(J_s)\chi$ and
\[
\textstyle\sum_k\de_k(b_s\xi_k)=\det(J_s)\sum_i(\de_i\tilde a)(\psi_s)\chi_i+(\tilde a\circ\psi_s)\bigl[\tr(\operatorname{adj}(J_s)D\chi)+\sum_i\chi_i\sum_k\de_k\operatorname{adj}(J_s)_{ki}\bigr],
\]
where the first term is $\det(J_s)\,\de_s(\tilde a\circ\psi_s)$, the trace is $\de_s\det J_s$ (Jacobi's formula, $\de_sJ_s=D\chi$) and the last sum vanishes (\cref{lem:fsp-piola}). By (a) and (b), $\cN'(s)=\SPn_0[\tilde F\circ\psi_s;\sum_k(\de_kh_k+z^{-1}\de_k(\tilde F\circ\psi_s)h_k)]$ with $h_k:=b_s\xi_k$, which vanishes by \cref{lem:fsp-ibp} over $A[s]$. So $\cN(1)=\cN(0)$. The full form follows from the branch rule, the critical values being equal.
\end{proof}

\begin{lem}[The uniformising map]\genusstatus{AI-proposed, unchecked}
\label{lem:fsp-cfu}
Let $\Bbbk$ be $\bbQ$ or $\bbQ(\zeta)$, $A:=\Bbbk\fps{u}$, $\psi(T;u)\in1+u\bbQ[T]\fps{u}$, $Y(X):=X\psi(X^5;u)$, $X_0\in\Bbbk$ with $X_0^5=1$, $Y_0:=Y(X_0)$, $f\in A[Y,Y^{-1}]$ and $\hat f:=f\circ Y$. Then $Y'(X_0)\in1+uA$, $X_0$ is a non-degenerate critical point of $\hat f$ over $A$ if and only if $Y_0$ is one of $f$, and then $\hat f(X_0)=f(Y_0)$ and, for every amplitude $\beta\in A\fps{Y-Y_0}\lser{\zh}$ and with $\zh$ in the role of $z$,
\[
\SPn_{X_0}\bigl[-\hat f;(\beta\circ Y)Y'/Y'(X_0)\bigr]=\SPn_{Y_0}[-f;\beta],\qquad\mathrm{SP}_{X_0}\bigl[\re^{-\hat f/\zh}(\beta\circ Y)Y'\,\dd X\bigr]=\mathrm{SP}_{Y_0}\bigl[\re^{-f/\zh}\beta\,\dd Y\bigr]
\]
in $A\lser{\zh}$, i.e.\ termwise in $u$, with the branch rule of \cref{lem:fsp-cov}.
\end{lem}

\begin{proof}
At each order in $u$, $Y$ is a polynomial in $X$, so $Y(X_0+t)-Y_0\in tA\fps{t}$; and $Y'(X_0)=\psi(1;u)+5\de_T\psi(1;u)\in1+uA$. As $Y_0\in X_0+uA$ is a unit, $Y(X_0+t)^{-1}\in A\fps{t}$, so $\hat f(X_0+t)\in A\fps{t}$. Since $\hat f'(X_0)=f'(Y_0)Y'(X_0)$, and $\hat f''(X_0)=f''(Y_0)Y'(X_0)^2$ when $f'(Y_0)=0$, the equivalence holds; the rest is \cref{lem:fsp-cov} over $A$ for the map $X\mapsto Y(X)$ and the phase $-f$ at $Y_0$.
\end{proof}

\begin{lem}[Iterated stationary phase]\genusstatus{AI-proposed, unchecked}
\label{lem:fsp-fubini}
Let $x=(v,w)$, $v=(v_1,\dots,v_{n_1})$, $w=(w_1,\dots,w_{n_2})$, let $F\in A\fps{x-p}$ have the non-degenerate critical point $p=(v_0,w_0)$ over $A$, with Hessian $\mathsf H=\bigl(\begin{smallmatrix}\mathsf H_{vv}&\mathsf H_{vw}\\ \mathsf H_{wv}&\mathsf H_{ww}\end{smallmatrix}\bigr)$ and $\mathsf H_{vv}\in GL_{n_1}(A)$, let $a\in A\fps{x-p}\lser{z}$ and $A_w:=A\fps{w-w_0}$.

(i) There is a unique $v_{\mathrm c}(w)$ in $v_0+((w-w_0)A_w)^{n_1}$ with $\de_vF(v_{\mathrm c}(w),w)=0$; the inner Hessian $\mathsf H_{vv}(w):=\de_v^2F(v_{\mathrm c}(w),w)$ lies in $GL_{n_1}(A_w)$, and $\mathsf H_{vv}(w_0)=\mathsf H_{vv}$.

(ii) $F^\flat(w):=F(v_{\mathrm c}(w),w)$ has $\nabla F^\flat(w_0)=0$ and $\Hess F^\flat(w_0)=\mathsf S:=\mathsf H_{ww}-\mathsf H_{wv}\mathsf H_{vv}^{-1}\mathsf H_{vw}$, and $\det\mathsf H=\det\mathsf H_{vv}\det\mathsf S$; hence $\det\mathsf S\in A^\times$, and $w_0$ is a non-degenerate critical point of $F^\flat$.

(iii) Let
\[
B(w;z):=\SPn_{v_{\mathrm c}(w)}[F(\cdot,w);a(\cdot,w)],\qquad\mathfrak b(w):=\bigl(\det\mathsf H_{vv}(w)/\det\mathsf H_{vv}\bigr)^{-1/2}
\]
be the inner series in $A_w\lser{z}$, computed over $A_w$, and the Fubini factor, a binomial series in $1+(w-w_0)A_w$. Then
\[
\SPn_p[F;a]=\SPn_{w_0}\bigl[F^\flat;\ \mathfrak b(w)\,B(w;z)\bigr];
\]
for full symbols,
\[
\mathrm{SP}_p\bigl[\re^{F/z}a\,\dd^{n_1}v\,\dd^{n_2}w\bigr]=\mathrm{SP}_{w_0}\bigl[\re^{F^\flat/z}(2\pi z)^{n_1/2}\det(-\mathsf H_{vv}(w))^{-1/2}B(w;z)\,\dd^{n_2}w\bigr],
\]
with $\det(-\mathsf H_{vv}(w))^{1/2}:=\det(-\mathsf H_{vv})^{1/2}\mathfrak b(w)^{-1}$ and $\det(-\mathsf H)^{1/2}:=\det(-\mathsf H_{vv})^{1/2}\det(-\mathsf S)^{1/2}$. So every splitting into an inner block with invertible Hessian and the rest computes the same total stationary phase; in particular two orders of iteration agree.
\end{lem}

\begin{proof}
(i) \Cref{lem:fsp-ift} over $A_w$, $I=(w-w_0)$, for $v'\mapsto F(v_0+v',w)$: the gradient at $v'=0$ vanishes at $w_0$, and the Hessian is $\mathsf H_{vv}$ modulo $I$.
(ii) $\nabla F^\flat(w)=\de_wF(v_{\mathrm c}(w),w)$, which vanishes at $w_0$; differentiating $\de_vF(v_{\mathrm c}(w),w)=0$ gives $Dv_{\mathrm c}(w_0)=-\mathsf H_{vv}^{-1}\mathsf H_{vw}$, so $\Hess F^\flat(w_0)=\mathsf H_{ww}+\mathsf H_{wv}Dv_{\mathrm c}(w_0)=\mathsf S$; and $\mathsf H=\bigl(\begin{smallmatrix}1&0\\ \mathsf H_{wv}\mathsf H_{vv}^{-1}&1\end{smallmatrix}\bigr)\diag(\mathsf H_{vv},\mathsf S)\bigl(\begin{smallmatrix}1&\mathsf H_{vv}^{-1}\mathsf H_{vw}\\0&1\end{smallmatrix}\bigr)$.
(iii) \emph{Straightening.} $\Theta(v',w):=(v_{\mathrm c}(w)+v',w)$ has $\det D\Theta=1$; by \cref{lem:fsp-cov} the left side is $\SPn_{(0,w_0)}[F\circ\Theta;a\circ\Theta]$, while $B$, $F^\flat$ and $\mathsf H_{vv}(w)$ are unchanged (the inner data are only translated). So we may assume $\de_vF(0,w)=0$ for all $w$, and $v_0=w_0=0$; then $\mathsf H_{vw}=0$, $\mathsf S=\mathsf H_{ww}$ and $F^\flat(w)=F(0,w)$.
\emph{Factorisation.} Expanding in $v$, $F(v,w)=F^\flat(w)+\frac12v^{\mathrm T}\mathsf H_{vv}(w)v+R(v,w)$, $R\in(v)^3$ over $A_w$. The total Hessian is $\diag(\mathsf H_{vv},\mathsf H_{ww})$ and the total cubic part is $F^\flat_{\geq3}(w)+\frac12v^{\mathrm T}\mathsf K(w)v+R$, with $F^\flat_{\geq3}:=F^\flat-F^\flat(0)-\frac12w^{\mathrm T}\mathsf H_{ww}w$ and $\mathsf K(w):=\mathsf H_{vv}(w)-\mathsf H_{vv}$, whose entries lie in $(w)$; each summand divided by $z$ has terms of weight $\geq1/2$. By \cref{lem:fsp-prod}, the linearity of $E^v$ over $\cR_{n_2}(A)$, the series in $w$ and $z$ (\cref{lem:fsp-ring}~(iv)), and \cref{lem:fsp-pg}~(iii) in $v$ with outer variables $w$,
\[
\SPn_0[F;a]=E^w_{-z\mathsf H_{ww}^{-1}}\Bigl[\re^{F^\flat_{\geq3}/z}\,E^v_{-z\mathsf H_{vv}^{-1}}\bigl[\re^{v^{\mathrm T}\mathsf Kv/(2z)}\re^{R/z}a\bigr]\Bigr]=E^w_{-z\mathsf H_{ww}^{-1}}\Bigl[\re^{F^\flat_{\geq3}/z}\,\mathfrak b(w)\,E^v_{-z\mathsf H_{vv}(w)^{-1}}\bigl[\re^{R/z}a\bigr]\Bigr].
\]
\emph{Rings.} At fixed $w$ the phase $v\mapsto F(v,w)-F^\flat(w)$ has Hessian $\mathsf H_{vv}(w)$ and cubic part $R(\cdot,w)$ at $v=0$, so $B(w;z)=E_{-z\mathsf H_{vv}(w)^{-1}}[\re^{R/z}a]$ in $\cR_{n_1}(A_w)$. The inner functional above is computed in the different ring $\cR_{n_1+n_2}(A)$; but $\re^{R/z}a$ is the same formal sum in both, both functionals act by the same rule on monomials, and each coefficient of $w^\beta z^k$ is the same finite sum. So the last expression is $\SPn_0[F^\flat;\mathfrak bB]$. The full form follows from the branch rules, as $F^\flat(w_0)=F(p)$ and $\det(-\mathsf H)=\det(-\mathsf H_{vv})\det(-\mathsf S)$.
\end{proof}

\input{sections/a_formalsp_5}



